\cleardoublepage
\selectlanguage{english}
\chapter*{Abstract}
\addcontentsline{toc}{chapter}{Abstract}
\label{front:abstract}

The analysis and qualification of public tenders requires the review of
extensive technical and administrative documentation in order to identify the
applicable requirements, assess the bidder's financial standing, and evaluate
the viability of the contract. When carried out manually, this process places
a significant burden on the responsible teams and may lead to errors of
omission.

This Bachelor's Thesis addresses the design and development of LIA, a
web-based commercial opportunity management platform focused on public
tenders. The solution organises opportunity portfolios through configurable
pipelines and structured workflows. It also integrates generative artificial
intelligence capabilities to support the extraction of key information, the
preliminary assessment of financial standing, and the generation of analysis
documents.

The overall objective of the project is to reduce the workload involved in
analysing public tenders and gathering their requirements. To this end, the
project aims to automate the assessment of financial standing, optimise the
cost-performance ratio of generative artificial intelligence providers, and
facilitate the creation of workflows adapted to tender analysis. The platform
has been developed as a Turborepo monorepository, with a Vue 3 web application
and a NestJS API. Technical validation combines unit, component, integration,
and end-to-end tests. The experimental evaluation compares 92 generations from
two Google models using quality, structural validity, latency, cost, and
hallucination criteria. The results support selecting Flash for summaries and
Flash-Lite for control questions, custom fields, and workflow decisions. The
deployment configuration retains Flash-Lite as its single global model by
developer decision, although the experiment recommends Flash for summaries. Outputs remain
subject to human review; evaluation on fictitious documents does not, by
itself, demonstrate a quantified reduction in analysis time in business use.

\vspace{0.5cm}
\noindent\textbf{Keywords:} public tenders, generative artificial
intelligence, opportunity management, financial standing assessment,
workflows.

\selectlanguage{spanish}
